%% bare_jrnl.tex
%% V1.4b
%% 2015/08/26
%% by Michael Shell
%% see http://www.michaelshell.org/
%% for current contact information.
%%
%% This is a skeleton file demonstrating the use of IEEEtran.cls
%% (requires IEEEtran.cls version 1.8b or later) with an IEEE
%% journal paper.
%%
%% Support sites:
%% http://www.michaelshell.org/tex/ieeetran/
%% http://www.ctan.org/pkg/ieeetran
%% and
%% http://www.ieee.org/

%%*************************************************************************
%% Legal Notice:
%% This code is offered as-is without any warranty either expressed or
%% implied; without even the implied warranty of MERCHANTABILITY or
%% FITNESS FOR A PARTICULAR PURPOSE! 
%% User assumes all risk.
%% In no event shall the IEEE or any contributor to this code be liable for
%% any damages or losses, including, but not limited to, incidental,
%% consequential, or any other damages, resulting from the use or misuse
%% of any information contained here.
%%
%% All comments are the opinions of their respective authors and are not
%% necessarily endorsed by the IEEE.
%%
%% This work is distributed under the LaTeX Project Public License (LPPL)
%% ( http://www.latex-project.org/ ) version 1.3, and may be freely used,
%% distributed and modified. A copy of the LPPL, version 1.3, is included
%% in the base LaTeX documentation of all distributions of LaTeX released
%% 2003/12/01 or later.
%% Retain all contribution notices and credits.
%% ** Modified files should be clearly indicated as such, including  **
%% ** renaming them and changing author support contact information. **
%%*************************************************************************


% *** Authors should verify (and, if needed, correct) their LaTeX system  ***
% *** with the testflow diagnostic prior to trusting their LaTeX platform ***
% *** with production work. The IEEE's font choices and paper sizes can   ***
% *** trigger bugs that do not appear when using other class files.       ***                          ***
% The testflow support page is at:
% http://www.michaelshell.org/tex/testflow/



\documentclass[journal]{IEEEtran}
%
% If IEEEtran.cls has not been installed into the LaTeX system files,
% manually specify the path to it like:
% \documentclass[journal]{../sty/IEEEtran}





% Some very useful LaTeX packages include:
% (uncomment the ones you want to load)


% *** MISC UTILITY PACKAGES ***
%
%\usepackage{ifpdf}
% Heiko Oberdiek's ifpdf.sty is very useful if you need conditional
% compilation based on whether the output is pdf or dvi.
% usage:
% \ifpdf
%   % pdf code
% \else
%   % dvi code
% \fi
% The latest version of ifpdf.sty can be obtained from:
% http://www.ctan.org/pkg/ifpdf
% Also, note that IEEEtran.cls V1.7 and later provides a builtin
% \ifCLASSINFOpdf conditional that works the same way.
% When switching from latex to pdflatex and vice-versa, the compiler may
% have to be run twice to clear warning/error messages.

% TODO: Remove xcolor package
% It was just added for highlighting fillers
\usepackage{xcolor}




% *** CITATION PACKAGES ***
%
%\usepackage{cite}
% cite.sty was written by Donald Arseneau
% V1.6 and later of IEEEtran pre-defines the format of the cite.sty package
% \cite{} output to follow that of the IEEE. Loading the cite package will
% result in citation numbers being automatically sorted and properly
% "compressed/ranged". e.g., [1], [9], [2], [7], [5], [6] without using
% cite.sty will become [1], [2], [5]--[7], [9] using cite.sty. cite.sty's
% \cite will automatically add leading space, if needed. Use cite.sty's
% noadjust option (cite.sty V3.8 and later) if you want to turn this off
% such as if a citation ever needs to be enclosed in parenthesis.
% cite.sty is already installed on most LaTeX systems. Be sure and use
% version 5.0 (2009-03-20) and later if using hyperref.sty.
% The latest version can be obtained at:
% http://www.ctan.org/pkg/cite
% The documentation is contained in the cite.sty file itself.






% *** GRAPHICS RELATED PACKAGES ***
%
\ifCLASSINFOpdf
  % \usepackage[pdftex]{graphicx}
  % declare the path(s) where your graphic files are
  % \graphicspath{{../pdf/}{../jpeg/}}
  % and their extensions so you won't have to specify these with
  % every instance of \includegraphics
  % \DeclareGraphicsExtensions{.pdf,.jpeg,.png}
\else
  % or other class option (dvipsone, dvipdf, if not using dvips). graphicx
  % will default to the driver specified in the system graphics.cfg if no
  % driver is specified.
  % \usepackage[dvips]{graphicx}
  % declare the path(s) where your graphic files are
  % \graphicspath{{../eps/}}
  % and their extensions so you won't have to specify these with
  % every instance of \includegraphics
  % \DeclareGraphicsExtensions{.eps}
\fi
% graphicx was written by David Carlisle and Sebastian Rahtz. It is
% required if you want graphics, photos, etc. graphicx.sty is already
% installed on most LaTeX systems. The latest version and documentation
% can be obtained at: 
% http://www.ctan.org/pkg/graphicx
% Another good source of documentation is "Using Imported Graphics in
% LaTeX2e" by Keith Reckdahl which can be found at:
% http://www.ctan.org/pkg/epslatex
%
% latex, and pdflatex in dvi mode, support graphics in encapsulated
% postscript (.eps) format. pdflatex in pdf mode supports graphics
% in .pdf, .jpeg, .png and .mps (metapost) formats. Users should ensure
% that all non-photo figures use a vector format (.eps, .pdf, .mps) and
% not a bitmapped formats (.jpeg, .png). The IEEE frowns on bitmapped formats
% which can result in "jaggedy"/blurry rendering of lines and letters as
% well as large increases in file sizes.
%
% You can find documentation about the pdfTeX application at:
% http://www.tug.org/applications/pdftex





% *** MATH PACKAGES ***
%
\usepackage{amsmath}
% A popular package from the American Mathematical Society that provides
% many useful and powerful commands for dealing with mathematics.
%
% Note that the amsmath package sets \interdisplaylinepenalty to 10000
% thus preventing page breaks from occurring within multiline equations. Use:
%\interdisplaylinepenalty=2500
% after loading amsmath to restore such page breaks as IEEEtran.cls normally
% does. amsmath.sty is already installed on most LaTeX systems. The latest
% version and documentation can be obtained at:
% http://www.ctan.org/pkg/amsmath





% *** SPECIALIZED LIST PACKAGES ***
%
%\usepackage{algorithmic}
% algorithmic.sty was written by Peter Williams and Rogerio Brito.
% This package provides an algorithmic environment fo describing algorithms.
% You can use the algorithmic environment in-text or within a figure
% environment to provide for a floating algorithm. Do NOT use the algorithm
% floating environment provided by algorithm.sty (by the same authors) or
% algorithm2e.sty (by Christophe Fiorio) as the IEEE does not use dedicated
% algorithm float types and packages that provide these will not provide
% correct IEEE style captions. The latest version and documentation of
% algorithmic.sty can be obtained at:
% http://www.ctan.org/pkg/algorithms
% Also of interest may be the (relatively newer and more customizable)
% algorithmicx.sty package by Szasz Janos:
% http://www.ctan.org/pkg/algorithmicx




% *** ALIGNMENT PACKAGES ***
%
%\usepackage{array}
% Frank Mittelbach's and David Carlisle's array.sty patches and improves
% the standard LaTeX2e array and tabular environments to provide better
% appearance and additional user controls. As the default LaTeX2e table
% generation code is lacking to the point of almost being broken with
% respect to the quality of the end results, all users are strongly
% advised to use an enhanced (at the very least that provided by array.sty)
% set of table tools. array.sty is already installed on most systems. The
% latest version and documentation can be obtained at:
% http://www.ctan.org/pkg/array


% IEEEtran contains the IEEEeqnarray family of commands that can be used to
% generate multiline equations as well as matrices, tables, etc., of high
% quality.




% *** SUBFIGURE PACKAGES ***
%\ifCLASSOPTIONcompsoc
%  \usepackage[caption=false,font=normalsize,labelfont=sf,textfont=sf]{subfig}
%\else
%  \usepackage[caption=false,font=footnotesize]{subfig}
%\fi
% subfig.sty, written by Steven Douglas Cochran, is the modern replacement
% for subfigure.sty, the latter of which is no longer maintained and is
% incompatible with some LaTeX packages including fixltx2e. However,
% subfig.sty requires and automatically loads Axel Sommerfeldt's caption.sty
% which will override IEEEtran.cls' handling of captions and this will result
% in non-IEEE style figure/table captions. To prevent this problem, be sure
% and invoke subfig.sty's "caption=false" package option (available since
% subfig.sty version 1.3, 2005/06/28) as this is will preserve IEEEtran.cls
% handling of captions.
% Note that the Computer Society format requires a larger sans serif font
% than the serif footnote size font used in traditional IEEE formatting
% and thus the need to invoke different subfig.sty package options depending
% on whether compsoc mode has been enabled.
%
% The latest version and documentation of subfig.sty can be obtained at:
% http://www.ctan.org/pkg/subfig




% *** FLOAT PACKAGES ***
%
%\usepackage{fixltx2e}
% fixltx2e, the successor to the earlier fix2col.sty, was written by
% Frank Mittelbach and David Carlisle. This package corrects a few problems
% in the LaTeX2e kernel, the most notable of which is that in current
% LaTeX2e releases, the ordering of single and double column floats is not
% guaranteed to be preserved. Thus, an unpatched LaTeX2e can allow a
% single column figure to be placed prior to an earlier double column
% figure.
% Be aware that LaTeX2e kernels dated 2015 and later have fixltx2e.sty's
% corrections already built into the system in which case a warning will
% be issued if an attempt is made to load fixltx2e.sty as it is no longer
% needed.
% The latest version and documentation can be found at:
% http://www.ctan.org/pkg/fixltx2e


%\usepackage{stfloats}
% stfloats.sty was written by Sigitas Tolusis. This package gives LaTeX2e
% the ability to do double column floats at the bottom of the page as well
% as the top. (e.g., "\begin{figure*}[!b]" is not normally possible in
% LaTeX2e). It also provides a command:
%\fnbelowfloat
% to enable the placement of footnotes below bottom floats (the standard
% LaTeX2e kernel puts them above bottom floats). This is an invasive package
% which rewrites many portions of the LaTeX2e float routines. It may not work
% with other packages that modify the LaTeX2e float routines. The latest
% version and documentation can be obtained at:
% http://www.ctan.org/pkg/stfloats
% Do not use the stfloats baselinefloat ability as the IEEE does not allow
% \baselineskip to stretch. Authors submitting work to the IEEE should note
% that the IEEE rarely uses double column equations and that authors should try
% to avoid such use. Do not be tempted to use the cuted.sty or midfloat.sty
% packages (also by Sigitas Tolusis) as the IEEE does not format its papers in
% such ways.
% Do not attempt to use stfloats with fixltx2e as they are incompatible.
% Instead, use Morten Hogholm'a dblfloatfix which combines the features
% of both fixltx2e and stfloats:
%
% \usepackage{dblfloatfix}
% The latest version can be found at:
% http://www.ctan.org/pkg/dblfloatfix




%\ifCLASSOPTIONcaptionsoff
%  \usepackage[nomarkers]{endfloat}
% \let\MYoriglatexcaption\caption
% \renewcommand{\caption}[2][\relax]{\MYoriglatexcaption[#2]{#2}}
%\fi
% endfloat.sty was written by James Darrell McCauley, Jeff Goldberg and 
% Axel Sommerfeldt. This package may be useful when used in conjunction with 
% IEEEtran.cls'  captionsoff option. Some IEEE journals/societies require that
% submissions have lists of figures/tables at the end of the paper and that
% figures/tables without any captions are placed on a page by themselves at
% the end of the document. If needed, the draftcls IEEEtran class option or
% \CLASSINPUTbaselinestretch interface can be used to increase the line
% spacing as well. Be sure and use the nomarkers option of endfloat to
% prevent endfloat from "marking" where the figures would have been placed
% in the text. The two hack lines of code above are a slight modification of
% that suggested by in the endfloat docs (section 8.4.1) to ensure that
% the full captions always appear in the list of figures/tables - even if
% the user used the short optional argument of \caption[]{}.
% IEEE papers do not typically make use of \caption[]'s optional argument,
% so this should not be an issue. A similar trick can be used to disable
% captions of packages such as subfig.sty that lack options to turn off
% the subcaptions:
% For subfig.sty:
% \let\MYorigsubfloat\subfloat
% \renewcommand{\subfloat}[2][\relax]{\MYorigsubfloat[]{#2}}
% However, the above trick will not work if both optional arguments of
% the \subfloat command are used. Furthermore, there needs to be a
% description of each subfigure *somewhere* and endfloat does not add
% subfigure captions to its list of figures. Thus, the best approach is to
% avoid the use of subfigure captions (many IEEE journals avoid them anyway)
% and instead reference/explain all the subfigures within the main caption.
% The latest version of endfloat.sty and its documentation can obtained at:
% http://www.ctan.org/pkg/endfloat
%
% The IEEEtran \ifCLASSOPTIONcaptionsoff conditional can also be used
% later in the document, say, to conditionally put the References on a 
% page by themselves.




% *** PDF, URL AND HYPERLINK PACKAGES ***
%
\usepackage{url}
% url.sty was written by Donald Arseneau. It provides better support for
% handling and breaking URLs. url.sty is already installed on most LaTeX
% systems. The latest version and documentation can be obtained at:
% http://www.ctan.org/pkg/url
% Basically, \url{my_url_here}.




% *** Do not adjust lengths that control margins, column widths, etc. ***
% *** Do not use packages that alter fonts (such as pslatex).         ***
% There should be no need to do such things with IEEEtran.cls V1.6 and later.
% (Unless specifically asked to do so by the journal or conference you plan
% to submit to, of course. )


% correct bad hyphenation here
\hyphenation{op-tical net-works semi-conduc-tor}


\begin{document}
%
% paper title
% Titles are generally capitalized except for words such as a, an, and, as,
% at, but, by, for, in, nor, of, on, or, the, to and up, which are usually
% not capitalized unless they are the first or last word of the title.
% Linebreaks \\ can be used within to get better formatting as desired.
% Do not put math or special symbols in the title.
\title{Our title \\ maybe 2 lines}
%
%
% author names and IEEE memberships
% note positions of commas and nonbreaking spaces ( ~ ) LaTeX will not break
% a structure at a ~ so this keeps an author's name from being broken across
% two lines.
% use \thanks{} to gain access to the first footnote area
% a separate \thanks must be used for each paragraph as LaTeX2e's \thanks
% was not built to handle multiple paragraphs
%

\author{Michael~Shell,~\IEEEmembership{Member,~IEEE,}
        John~Doe,~\IEEEmembership{Fellow,~OSA,}
        and~Jane~Doe,~\IEEEmembership{Life~Fellow,~IEEE}% <-this % stops a space
\thanks{M. Shell was with the Department
of Electrical and Computer Engineering, Georgia Institute of Technology, Atlanta,
GA, 30332 USA e-mail: (see http://www.michaelshell.org/contact.html).}% <-this % stops a space
\thanks{J. Doe and J. Doe are with Anonymous University.}% <-this % stops a space
\thanks{Manuscript received April 19, 2005; revised August 26, 2015.}}

% note the % following the last \IEEEmembership and also \thanks - 
% these prevent an unwanted space from occurring between the last author name
% and the end of the author line. i.e., if you had this:
% 
% \author{....lastname \thanks{...} \thanks{...} }
%                     ^------------^------------^----Do not want these spaces!
%
% a space would be appended to the last name and could cause every name on that
% line to be shifted left slightly. This is one of those "LaTeX things". For
% instance, "\textbf{A} \textbf{B}" will typeset as "A B" not "AB". To get
% "AB" then you have to do: "\textbf{A}\textbf{B}"
% \thanks is no different in this regard, so shield the last } of each \thanks
% that ends a line with a % and do not let a space in before the next \thanks.
% Spaces after \IEEEmembership other than the last one are OK (and needed) as
% you are supposed to have spaces between the names. For what it is worth,
% this is a minor point as most people would not even notice if the said evil
% space somehow managed to creep in.



% The paper headers
\markboth{EAITHD26 (Enhancing AI Through Hybrid Datasets),~Vol.~1, No.~1, September~2026}%
{Shell \MakeLowercase{\textit{et al.}}: Bare Demo of IEEEtran.cls for IEEE Journals}
% The only time the second header will appear is for the odd numbered pages
% after the title page when using the twoside option.
% 
% *** Note that you probably will NOT want to include the author's ***
% *** name in the headers of peer review papers.                   ***
% You can use \ifCLASSOPTIONpeerreview for conditional compilation here if
% you desire.




% If you want to put a publisher's ID mark on the page you can do it like
% this:
%\IEEEpubid{0000--0000/00\$00.00~\copyright~2015 IEEE}
% Remember, if you use this you must call \IEEEpubidadjcol in the second
% column for its text to clear the IEEEpubid mark.



% use for special paper notices
%\IEEEspecialpapernotice{(Invited Paper)}




% make the title area
\maketitle

% As a general rule, do not put math, special symbols or citations
% in the abstract or keywords.
\begin{abstract}
The abstract goes here.
\end{abstract}

% Note that keywords are not normally used for peerreview papers.
\begin{IEEEkeywords}
IEEE, IEEEtran, journal, \LaTeX, paper, template.
\end{IEEEkeywords}






% For peer review papers, you can put extra information on the cover
% page as needed:
% \ifCLASSOPTIONpeerreview
% \begin{center} \bfseries EDICS Category: 3-BBND \end{center}
% \fi
%
% For peerreview papers, this IEEEtran command inserts a page break and
% creates the second title. It will be ignored for other modes.
\IEEEpeerreviewmaketitle



\section{Introduction}
% The very first letter is a 2 line initial drop letter followed
% by the rest of the first word in caps.
% 
% form to use if the first word consists of a single letter:
% \IEEEPARstart{A}{demo} file is ....
% 
% form to use if you need the single drop letter followed by
% normal text (unknown if ever used by the IEEE):
% \IEEEPARstart{A}{}demo file is ....
% 
% Some journals put the first two words in caps:
% \IEEEPARstart{T}{his demo} file is ....
% 
% Here we have the typical use of a "T" for an initial drop letter
% and "HIS" in caps to complete the first word.
\IEEEPARstart{T}{his} \textcolor{gray}{introduction still needs to be written. Bla bla. CATCHING INTRO - maybe just shift II. RESEARCH QUESTION PART to intro? Bla Bla. STUDY APPROACH}
% You must have at least 2 lines in the paragraph with the drop letter
% (should never be an issue)
Ten shift quantification metrics were applied to nine hybrid dataset configurations for daytime urban street scene segmentation. Correlation with downstream performance was measured across two model architectures with three random states each. \textcolor{gray}{Bla Bla MAIN FINDINGS}


% An example of a floating figure using the graphicx package.
% Note that \label must occur AFTER (or within) \caption.
% For figures, \caption should occur after the \includegraphics.
% Note that IEEEtran v1.7 and later has special internal code that
% is designed to preserve the operation of \label within \caption
% even when the captionsoff option is in effect. However, because
% of issues like this, it may be the safest practice to put all your
% \label just after \caption rather than within \caption{}.
%
% Reminder: the "draftcls" or "draftclsnofoot", not "draft", class
% option should be used if it is desired that the figures are to be
% displayed while in draft mode.
%
%\begin{figure}[!t]
%\centering
%\includegraphics[width=2.5in]{myfigure}
% where an .eps filename suffix will be assumed under latex, 
% and a .pdf suffix will be assumed for pdflatex; or what has been declared
% via \DeclareGraphicsExtensions.
%\caption{Simulation results for the network.}
%\label{fig_sim}
%\end{figure}

% Note that the IEEE typically puts floats only at the top, even when this
% results in a large percentage of a column being occupied by floats.


% An example of a double column floating figure using two subfigures.
% (The subfig.sty package must be loaded for this to work.)
% The subfigure \label commands are set within each subfloat command,
% and the \label for the overall figure must come after \caption.
% \hfil is used as a separator to get equal spacing.
% Watch out that the combined width of all the subfigures on a 
% line do not exceed the text width or a line break will occur.
%
%\begin{figure*}[!t]
%\centering
%\subfloat[Case I]{\includegraphics[width=2.5in]{box}%
%\label{fig_first_case}}
%\hfil
%\subfloat[Case II]{\includegraphics[width=2.5in]{box}%
%\label{fig_second_case}}
%\caption{Simulation results for the network.}
%\label{fig_sim}
%\end{figure*}
%
% Note that often IEEE papers with subfigures do not employ subfigure
% captions (using the optional argument to \subfloat[]), but instead will
% reference/describe all of them (a), (b), etc., within the main caption.
% Be aware that for subfig.sty to generate the (a), (b), etc., subfigure
% labels, the optional argument to \subfloat must be present. If a
% subcaption is not desired, just leave its contents blank,
% e.g., \subfloat[].


% An example of a floating table. Note that, for IEEE style tables, the
% \caption command should come BEFORE the table and, given that table
% captions serve much like titles, are usually capitalized except for words
% such as a, an, and, as, at, but, by, for, in, nor, of, on, or, the, to
% and up, which are usually not capitalized unless they are the first or
% last word of the caption. Table text will default to \footnotesize as
% the IEEE normally uses this smaller font for tables.
% The \label must come after \caption as always.
%
%\begin{table}[!t]
%% increase table row spacing, adjust to taste
%\renewcommand{\arraystretch}{1.3}
% if using array.sty, it might be a good idea to tweak the value of
% \extrarowheight as needed to properly center the text within the cells
%\caption{An Example of a Table}
%\label{table_example}
%\centering
%% Some packages, such as MDW tools, offer better commands for making tables
%% than the plain LaTeX2e tabular which is used here.
%\begin{tabular}{|c||c|}
%\hline
%One & Two\\
%\hline
%Three & Four\\
%\hline
%\end{tabular}
%\end{table}


% Note that the IEEE does not put floats in the very first column
% - or typically anywhere on the first page for that matter. Also,
% in-text middle ("here") positioning is typically not used, but it
% is allowed and encouraged for Computer Society conferences (but
% not Computer Society journals). Most IEEE journals/conferences use
% top floats exclusively. 
% Note that, LaTeX2e, unlike IEEE journals/conferences, places
% footnotes above bottom floats. This can be corrected via the
% \fnbelowfloat command of the stfloats package.

\section{Research question}
Hybrid datasets, which combine real and synthetic data, are often used to address data-associated limitations in model performance and generalization. However, selecting an effective hybrid dataset configuration is costly when based on actual downstream performance, because it requires training a separate model for each candidate configuration. A more efficient alternative is to predict downstream performance. This study focuses on \textbf{distribution-shift-based prediction} with the following research question: 

\textit{Can pre-training shift quantification predict the utility of a hybrid dataset reliably and therefore replace naive selection among candidate configurations?}

\section{Methods}
This study evaluates the proposed approach in the context of \textbf{daytime urban street scene segmentation}.

\subsection{Data preparation}
\textcolor{gray}{Bla Bla: Describe the datasets (Cityscapes, Synscapes, GTA-V)} 

\textcolor{gray}{Bla Bla: Describe the preparation (number of samples selected, how samples were selected, resizing to uniform resolution, construction of the 9 hybrid train datasets + val + test datasets and the reasoning behind it, certain class indices are ignored - 19 instead of 34 classes)} 

\subsection{Shift quantification}
A distribution shift occurs when two data distributions differ. In machine learning, this typically refers to the shift between the distribution a model is trained on (training distribution) and the distribution it will be applied to (inference distribution, approximated by the test dataset). Hybrid datasets are often motivated by the goal of reducing this distribution shift: Adding synthetic data to real data can enrich the training distribution in a controlled manner, improving coverage and diversity, and supporting generalization to shifted inference conditions. Conversely, adding real data to synthetic data can help anchoring the model in real-world statistics, reducing the reality gap and mitigating reliance on unrealistic or harmful artifacts. 

Beyond this general notion, a distribution shift can be decomposed into several distinct components. Following the formulation of \cite{liu2022}, a segmentation model can be expressed as $G(x)=C(F(x))$, where the classifier $C$ outputs the posterior probability $P(Y|F(x))$ for an input image $x$. Applying Bayes' theorem yields 
\begin{equation}
\label{eq:bayes_shift}
P(Y \mid F(x)) = \frac{P(F(x) \mid Y) \cdot P(Y)}{P(F(x))}.
\end{equation}
The Bayes decomposition exposes the distributions that may differ between datasets. Based on the affected distribution, we can differentiate five shift types, of which this study focuses on three: covariate shift, label shift, and representation shift.\footnote{Additional shift types include class-conditional representation shift, which affects the input representation for a specific class $P(F(x)|Y)$, and concept shift, which affects the labeling function $P(Y|F(x))$ or $P(Y|X)$. Class-conditional representation shift is excluded because this study quantifies the shift at the dataset level rather than per class. Concept shift is excluded because all datasets share the same semantic label definitions.} Based on the information used for quantification, we group shift quantification metrics into three categories, each capturing a different aspect of the shift:

\textbf{Image-based metrics} quantify the covariate shift, that is, differences in the input distribution $P(X)$ that arise from changes in visual properties such as color and texture. These metrics operate on the raw images. This study considers two metrics that focus on a shift in color ($\text{js}_{\text{color}}$) and texture ($\text{js}_{\text{texture}}$) respectively. Covariate shift was quantified using Jensen-Shannon distance \eqref{eq:js} between the normalized mean histograms $p_{\text{train}}$, $p_{\text{test}}$. Per-image color histograms were computed in HSV space. Per-image texture histograms were computed using multiscale Local Binary Patterns (LBP) at two scales ($P=8$, $R=1$) and ($P=16$, $R=2$).

\begin{equation}
\label{eq:js}
\begin{split}
d_{\text{JS}}(p_{\text{train}}, p_{\text{test}})
= \sqrt{
\frac{1}{2} D_{\text{KL}}\!\left( p_{\text{train}} \,\middle\|\, m \right)
+
\frac{1}{2} D_{\text{KL}}\!\left( p_{\text{test}} \,\middle\|\, m \right)
}, \\
\text{with } m = \frac{1}{2}\left(p_{\text{train}} + p_{\text{test}}\right)
\end{split}
\end{equation}

\textbf{Label-based metrics} quantify the label shift, that is, differences in the class-label distribution $P(Y)$. These metrics operate on the segmentation masks. This study considers three metrics that focus on a shift in class frequency ($\text{js}_{\text{classFrequency}}$), class presence ($\text{js}_{\text{classPresence}}$), and scene complexity ($\text{js}_{\text{sceneComplexity}}$) respectively. Label shift was quantified using Jenson-Shannon distance \eqref{eq:js} between the normalized histograms $p_{\text{train}}$, $p_{\text{test}}$. The class frequency histogram of a dataset shows the fraction of all pixels belonging to each class. The class presence histogram of a dataset shows the fraction of images containing each class. The scene complexity histogram of a dataset shows the distribution of distinct-class counts per image.

\textbf{Representation-based metrics} quantify the representation shift, that is, differences in the feature distribution $P(F(x))$. These metrics operate in a learned feature space on image representations produced by a feature extractor $F$. This study used two pretrained feature extractors: DINOv2‑Base \cite{oquab2024} for appearance‑focused representations and StreetCLIP \cite{haas2023} for semantic‑focused representations. They were selected because both were exposed to street-view-like imagery during training, making them well aligned with this study's target domain. Representation shift was quantified using centroid distance \eqref{eq:centroid} and Fréchet distance \eqref{eq:frechet} \cite{NIPS2017_8a1d6947} between two distributions with means $\mu_{\text{train}}$, $\mu_{\text{test}}$ and covariance matrices $\Sigma_{\text{train}}$, $\Sigma_{\text{test}}$, computed in DINOv2-Base and StreetCLIP feature spaces. This yielded four metric variants: $\text{fréchet}_{\text{dinov2}}$, $\text{centroid}_{\text{dinov2}}$, $\text{fréchet}_{\text{streetclip}}$, and $\text{centroid}_{\text{streetclip}}$.

\begin{equation}
\label{eq:centroid}
d_{\text{centroid}}(\mu_{\text{train}},\mu_{\text{test}}) = \|\mu_{\text{train}} - \mu_{\text{test}}\|
\end{equation}

\begin{equation}
\label{eq:frechet}
\begin{aligned}
&d_{\text{Fréchet}}(\mu_{\text{train}},\Sigma_{\text{train}};\mu_{\text{test}},\Sigma_{\text{test}}) = \\
&\sqrt{
\|\mu_{\text{train}} - \mu_{\text{test}}\|^2
+ \mathrm{Tr}\!\left(\Sigma_{\text{train}} + \Sigma_{\text{test}} - 2(\Sigma_{\text{train}} \Sigma_{\text{test}})^{\frac{1}{2}}\right)
}
\end{aligned}
\end{equation}

\smallskip
The last representation-based metric considered in this study is \textbf{SADGE} (Structural and Appearance Domain Gap Estimator), which was shown to predict downstream model performance most reliably among the shift quantification metrics tested in \cite{bartkowiak2026}. SADGE differs fundamentally from the nine other metrics evaluated here: While all other metrics quantify distributional distance between datasets (by comparing dataset-level statistics in image, label, or feature space), SADGE computes pairwise similarity between individual images and aggregates these similarities into a dataset-level score. In our adaptation, each test image is compared to ten randomly selected candidate images from the hybrid training dataset. The strongest appearance and geometry match per test image is retained. Appearance similarity is measured using cosine similarity between DINOv3 \cite{simeoni2025dinov3} embeddings. Geometric similarity is measured using MASt3R \cite{mast3r} inlier counts. These per-image scores are averaged (to obtain dataset-level appearance similarity $\bar{A}$ and geometry similarity $\bar{G}$), standardized using z-score normalization \eqref{eq:zscore}, and fused through the pre-calibrated linear-bilinear model \eqref{eq:sadge} provided in \cite{bartkowiak2026}.

\begin{equation} \label{eq:zscore}
\begin{aligned}
\hat{A} &= \frac{\bar{A} - \mu_A}{\sigma_A}, && \text{with } \mu_A = 0.6359, \, \sigma_A = 0.1918 \\
\hat{G} &= \frac{\log(1+\bar{G}) - \mu_G}{\sigma_G}, && \text{with } \mu_G = 7.9420, \, \sigma_G = 1.7384
\end{aligned}
\end{equation}

\begin{equation}
\label{eq:sadge}
\begin{split}
\text{SADGE}(\text{ds}_\text{test}, \text{ds}_\text{train}) = a \cdot \hat{G} + b \cdot \hat{A} + c \cdot \hat{G} \cdot \hat{A}
\\
\text{with } a=0.0\text{, } b=1.8548\text{, } c=1.3399
\end{split}
\end{equation}

\subsection{Model training}
\textcolor{gray}{Bla Bla: Describe the segmentation models} 

\textcolor{gray}{Bla Bla: Describe the model training (hyperparameter settings, callbacks, loss function, recorded metrics, etc.)} 

\textcolor{gray}{Bla Bla: Describe the evaluation of model performance at inference time, e.g. metrics recorded etc.} 

\section{Results}
\textcolor{gray}{Bla Bla: Describe the observed general dataset differences between Cityscapes, Synscapses, GTA-V - maybe plot histograms or sth like that }

\textcolor{gray}{Bla Bla: Plot the training, validation loss progression to show that the models reach a stable state and to compare hybrid dataset configs already based on training progression}

\textcolor{gray}{Bla Bla: Main correlation analysis - plot shift scores against model performance scores and measure Pearson/Spearman correlation - identify promising metrics - do promising metrics differ between synthetic datasets? - are the results stable across seeds and model architectures?}

\section{Evaluation}
\textcolor{gray}{Bla Bla: Answer research question based on findings}

\textcolor{gray}{Bla Bla: Relate findings to previous research}

\textcolor{gray}{Bla Bla: Discuss limitations of the study}

\section{Conclusion}
\textcolor{gray}{Bla Bla: Conclusion/Summary}

\section{Reproducibility}
\textcolor{gray}{Bla Bla: Link GitHub repo, describe it shortly, give broad instructions on how to replicate our study}




% if have a single appendix:
%\appendix[Proof of the Zonklar Equations]
% or
%\appendix  % for no appendix heading
% do not use \section anymore after \appendix, only \section*
% is possibly needed

% use appendices with more than one appendix
% then use \section to start each appendix
% you must declare a \section before using any
% \subsection or using \label (\appendices by itself
% starts a section numbered zero.)
%


\appendices
\section{Related research}
Appendix one text goes here.

% you can choose not to have a title for an appendix
% if you want by leaving the argument blank
\section{}
Appendix two text goes here.


% use section* for acknowledgment
\section*{Acknowledgment}


The authors would like to thank...


% Can use something like this to put references on a page
% by themselves when using endfloat and the captionsoff option.
\ifCLASSOPTIONcaptionsoff
  \newpage
\fi



% trigger a \newpage just before the given reference
% number - used to balance the columns on the last page
% adjust value as needed - may need to be readjusted if
% the document is modified later
%\IEEEtriggeratref{8}
% The "triggered" command can be changed if desired:
%\IEEEtriggercmd{\enlargethispage{-5in}}

% REFERENCES section
% define references style (use offical IEEE format)
\bibliographystyle{IEEEtran}
% connect references.bib file
\bibliography{references}

% can use a bibliography generated by BibTeX as a .bbl file
% BibTeX documentation can be easily obtained at:
% http://mirror.ctan.org/biblio/bibtex/contrib/doc/
% The IEEEtran BibTeX style support page is at:
% http://www.michaelshell.org/tex/ieeetran/bibtex/
%\bibliographystyle{IEEEtran}
% argument is your BibTeX string definitions and bibliography database(s)
%\bibliography{IEEEabrv,../bib/paper}
%
% <OR> manually copy in the resultant .bbl file
% set second argument of \begin to the number of references
% (used to reserve space for the reference number labels box)
% \begin{thebibliography}{1}

% \bibitem{IEEEhowto:kopka}
% H.~Kopka and P.~W. Daly, \emph{A Guide to \LaTeX}, % 3rd~ed.\hskip 1em plus
%   0.5em minus 0.4em\relax Harlow, England: Addison-Wesley, 1999.

% \end{thebibliography}

% biography section
% 
% If you have an EPS/PDF photo (graphicx package needed) extra braces are
% needed around the contents of the optional argument to biography to prevent
% the LaTeX parser from getting confused when it sees the complicated
% \includegraphics command within an optional argument. (You could create
% your own custom macro containing the \includegraphics command to make things
% simpler here.)
%\begin{IEEEbiography}[{\includegraphics[width=1in,height=1.25in,clip,keepaspectratio]{mshell}}]{Michael Shell}
% or if you just want to reserve a space for a photo:

\begin{IEEEbiography}{Michael Shell}
Biography text here.
\end{IEEEbiography}

% if you will not have a photo at all:
\begin{IEEEbiographynophoto}{John Doe}
Biography text here.
\end{IEEEbiographynophoto}

% insert where needed to balance the two columns on the last page with
% biographies
%\newpage

\begin{IEEEbiographynophoto}{Jane Doe}
Biography text here.
\end{IEEEbiographynophoto}

% You can push biographies down or up by placing
% a \vfill before or after them. The appropriate
% use of \vfill depends on what kind of text is
% on the last page and whether or not the columns
% are being equalized.

%\vfill

% Can be used to pull up biographies so that the bottom of the last one
% is flush with the other column.
%\enlargethispage{-5in}



% that's all folks
\end{document}


